\section{Glossary}
\label{app:glossary}

{\footnotesize \setlength{\columnsep}{20pt}
\begin{multicols}{2}
\begin{description}[style=sameline, leftmargin=0pt, itemsep=1.5pt, font=\normalfont\bfseries]
\item[Ablation.] Replacing the value of a node by zero, by its mean or by its value on another prompt, and
measuring the change of the output (\cref{sec:interventions}).
\item[Activation.] Any intermediate value that the model computes on a given input.
\item[Activation patching.] Running the model on one prompt while overwriting one node with its value from a
second, minimally different prompt (\cref{sec:patching}).
\item[Attention head.] The component that moves information between positions. It is described by the matrices
$W_{QK}$ (from where to read) and $W_{OV}$ (what to write), see \eqref{eq:qkov}.
\item[Attention pattern.] The weights $A^{\ell,h}_{ij}$ with which head $h$ of layer $\ell$ at destination $i$
reads from the sources $j\le i$.
\item[Attribution graph.] A graph of transcoder latents and their direct effects on one another for a single
prompt (\cref{sec:featurecircuits}).
\item[Attribution patching.] The first-order estimate \eqref{eq:attrpatch} of the effect of activation patching.
\item[Circuit.] A subgraph of the model's components together with a description of what each one computes on
a task (\cref{def:circuit}).
\item[Clean and corrupted prompt.] The two prompts of a patching experiment. They differ only in the aspect
under study.
\item[Composition.] The use of the output of an earlier head in the query, key or value of a later head.
\item[Dictionary direction.] A column $\Wdec[j]$ of the decoder of a sparse autoencoder.
\item[Direct logit attribution.] The term that one component contributes to a logit through its own addition
to the residual stream (\cref{prop:dla}).
\item[Faithful.] Said of an explanation that correctly describes what the model computes. Opposed to merely
plausible.
\item[Feature.] A property of the input or of the intermediate computation that the model represents.
\item[Hook.] A place in the code of a model at which an activation can be read or overwritten.
\item[Latent.] One unit of a sparse autoencoder or transcoder. A candidate for a feature.
\item[Linear probe.] A logistic regression that predicts a label from an activation (\cref{sec:probing}).
\item[Logit.] The score of a token before the softmax \eqref{eq:softmax}.
\item[Logit difference.] The logit of the correct token minus that of a chosen alternative. Our metric
\eqref{eq:ld}.
\item[Logit lens.] Applying the final read-out to an intermediate residual stream \eqref{eq:lens}.
\item[Neuron.] One coordinate of the hidden layer of an MLP.
\item[Polysemantic.] Said of a neuron that responds to several unrelated features.
\item[Privileged basis.] A coordinate system that the architecture singles out, as the neurons of an MLP. The
residual stream has none.
\item[Residual stream.] The vector $x_i^{(\ell)}$ that is kept for each position and to which all components
add their output.
\item[Sparse autoencoder.] A wide one-layer autoencoder trained to reconstruct activations with few active
latents (\cref{sec:sae}).
\item[Steering.] Adding a direction to the residual stream, or removing it, to change a behaviour
(\cref{sec:steering}).
\item[Superposition.] The storage of more features than dimensions along almost orthogonal directions
(\cref{sec:superposition}).
\item[Token.] A word or word fragment. The unit in which the model reads and writes text.
\item[Transcoder.] A sparse, wide replacement of an MLP that predicts its output from its input.
\item[Unembedding.] The matrix $W_U$ that turns the final residual stream into logits.
\end{description}
\end{multicols}
}

\section{Solutions to the exercises}
\label{app:solutions}

{\small
\paragraph{\Cref{ex:params}.} A head has $4\cdot d\cdot d_{\mathrm{head}}=4\cdot768\cdot64=196{,}608$ weights. $W_E$ has
$50257\cdot768=38{,}597{,}376$ entries, which is $31\%$ of the model.

\paragraph{\Cref{ex:ones}.} Since $\mean(x+s\,\one)=\mean(x)+s$, we get
$x+s\,\one-\mean(x+s\,\one)\,\one=x-\mean(x)\,\one$. The numerator and the scale in \eqref{eq:ln} are therefore
unchanged. Every component reads the stream through $\LN$, so all component outputs stay the same. By
\eqref{eq:update}, all later streams at this position differ from the original ones by the same $s\,\one$, which
the final $\LN$ in \eqref{eq:unembed} removes.

\paragraph{\Cref{ex:rank}.} $W_{OV}=W_OW_V$ factors through $\R^{d_{\mathrm{head}}}$, so its rank is at most
$d_{\mathrm{head}}=64$. Whatever a head writes lies in the column space of $W_O$, a fixed subspace of dimension at
most $64$ of $\R^{768}$.

\paragraph{\Cref{ex:linearig}.} The gradient $\nabla_{z_i}F=w_i$ is constant, so
$\mathrm{IG}_i(z)=z_i^\top\int_0^1w_i\,du=w_i^\top z_i$, which is gradient times input. For $F(z)=z^2$ we have
$F'(z)=2z$. Gradient times input is $2z^2$, and $\mathrm{IG}(z)=z\int_0^12uz\,du=z^2=F(z)-F(0)$.

\paragraph{\Cref{ex:lenslast}.} At $\ell=L$, \eqref{eq:lens} is \eqref{eq:unembed}. The components up to layer 8
add up to $x_n^{(9)}$ by the proof of \cref{prop:sum}. Their DLA terms therefore add up to
$(W_U[\mathrm{IO}]-W_U[\mathrm{S}])^\top(x_n^{(9)}-\mean(x_n^{(9)})\,\one)/\scale(x_n^{(L)})$. The lens divides the same
numerator by $\scale(x_n^{(9)})$. The two numbers differ by the factor $\scale(x_n^{(9)})/\scale(x_n^{(L)})$.

\paragraph{\Cref{ex:probecontrol}.} Being a first name is a function of the current token, and the token is
encoded in $x_i^{(0)}$ by \eqref{eq:embed}. Train the same probe on layer $0$, and on a control task that gives
each token type a random label. The probe on layer $0$ will succeed as well, so the result shows nothing
specific about layer 6.

\paragraph{\Cref{ex:triangle}.} $v_i^\top v_j=\cos120^\circ=-1/2$. If only feature $i$ is active, $x=f_iv_i$. Then
$\relu(v_i^\top x)=f_i$ and, for $j\ne i$, $\relu(v_j^\top x)=\relu(-f_i/2)=0$. If no feature is active, $x=0$.

\paragraph{\Cref{ex:bound}.} $N\le e^{10.24}\approx2.8\times10^4$ for $\varepsilon=0.1$ and
$N\le e^{40.96}\approx6\times10^{17}$ for $\varepsilon=0.2$. The condition $k\varepsilon<1/2$ allows $k\le4$ and
$k\le2$.

\paragraph{\Cref{ex:nobias}.} Now $\hat f_2=0$ and the loss is $\E f_2^2=\delta/3$, which exceeds $\delta^2/3$ for
every $\delta<1$. The bias lets the model output the mean of the feature it ignores. This lowers the error from
the second moment $\delta/3$ to the variance $\delta/3-\delta^2/4$.

\paragraph{\Cref{ex:normconstraint}.} Since $\relu(y/2)=\relu(y)/2$, the new activation is $f_j(x)/2$. Its product
with the doubled column is unchanged, and its penalty is halved. Without the constraint, training would let the
columns grow and the activations shrink towards zero, and the penalty would lose its meaning.

\paragraph{\Cref{ex:fvu}.} The optimal value is $f=\max(0,f^\star-0.2)$. The error is $f^\star$ for $f^\star<0.2$
and $0.2$ otherwise. Hence $\E\,(f^\star-f)^2=\int_0^{0.2}u^2\,du+0.8\cdot0.04=0.0027+0.032\approx0.035$.

\paragraph{\Cref{ex:ggb}.} Observations: weak and medium activations are also about the bridge, and the latent
fires on most mentions of the bridge in a new corpus. Intervention: amplifying the latent makes the model talk
about the bridge, and a random latent of equal strength does not.

\paragraph{\Cref{ex:andor}.} (a) Denoising gives $\max(1,0)=1$, effect $1$. Noising gives $\max(0,1)=1$, effect
$0$. Each node is sufficient and none is necessary. (b) Denoising gives $\min(1,0)=0$, effect $0$. Noising
gives $\min(0,1)=0$, effect $1$. Each node is necessary and none is sufficient.

\paragraph{\Cref{ex:affine}.} An affine function has the same gradient everywhere, and its increment between two
points equals this gradient times their difference. Example: the node $\LN(x_n^{(L)})$ with the logit difference
as metric, by \eqref{eq:ld}.

\paragraph{\Cref{ex:projection}.} The map is $x\mapsto(I-vv^\top)x$, which is linear.
$(I-vv^\top)^2=I-2vv^\top+v(v^\top v)v^\top=I-vv^\top$. And $w^\top(I-vv^\top)x=w^\top x-(w^\top v)(v^\top x)=w^\top x$
when $w^\top v=0$.

\paragraph{\Cref{ex:skiptrigram}.} Both \tok{in} and \tok{at} attend to \tok{keep}. What the head then writes
depends only on the source \tok{keep}, and it promotes both \tok{mind} and \tok{bay}. So \tok{bay} is also
promoted after \tok{in}. To avoid this, the table $W_U[t]^\top W_{OV}^{0,h}W_E[t_j]$ would have to depend on the
destination token.

\paragraph{\Cref{ex:induction}.} The induction head has to attend to the position $j$ whose \emph{previous} token
equals the current token. In one layer, the score of the pair $(i,j)$ depends on $t_i$, $t_j$ and the positions,
and not on $t_{j-1}$. An earlier layer must first write $t_{j-1}$ into position $j$.

\paragraph{\Cref{ex:cosprime}.} $\cos(2\pi\cdot14q/113)=1$ holds exactly when $113$ divides $14q$, which means
that $113$ divides $q$, since $113$ is prime. For $P=112$ the number $14q/112=q/8$ is an integer for $q=8$, so a
wrong answer would score as high as the correct one at this frequency.

\paragraph{\Cref{ex:critique}.} A control task or shuffled labels, and a check that lying and honest prompts do
not differ in topic or length (modes 1 and 2). An intervention along the probe direction that changes whether
the model lies, compared with a random direction (modes 1 and 6). A test on lies of a different kind than those
in the training set (mode 8).
}
